\chapter*{Conclusion Générale}
\addcontentsline{toc}{chapter}{Conclusion Générale}

Le présent rapport de Projet de Fin d'Études rend compte de la conceptualisation, conception et implémentation rigoureuses de Guidni, une plateforme touristique novatrice cherchant à automatiser la complexité logistique des séjours adressés à l'île de Djerba. En dépassant très explicitement le simple paradigme d'une réservation transactionnelle de commerce classique, notre objectif scientifique central gravitait autour de la consolidation de l'asynchronie informatique via l'assemblage simultané d'un agent de dialogue doué de compréhension conversationnelle (grands modèles de langages) et d'un propulseur d'exploitation contextuelle mathématique d'apprentissage par l'expérience (bandits manchots).

Pour faire face à cette équation conceptuelle riche en risques matériels algorithmiques potentiels, l'adoption fondatrice de la structure méthodologique Agile par le framework Scrum s'est prouvée déterminante. Au cours des cinq bilans d'avancements programmés minutieusement sur l'étendue des mois, d'octobre à février, cette structuration a tempéré nos incertitudes grandissantes. Qu'il s'agissait du maillage des nœuds directifs LangGraph au profit d'une intelligence d'encadrement en dix-huit modules intégrés, de la conception minutieuse du modèle tensoriel asynchrone RAG multilinguistique dotés d'accélération sur fil logiques de réindexation, ou encore du codage pur évitant une latence relationnelle par ancrage mathématique LinUCB en RAM ; la définition fragmentée du produit selon le critère de réussite immuable a guidé inlassablement la mise en production.

Les contributions techniques portées soutiennent une avancée significative sur la manipulation de composants statistiques en charge au sein d'orchestrateurs backend Python sous la coupole FastAPI. Assurément, la méthode de calcul du score par croisement du profil de type d'usager et de l'historique de clics encadrant les descripteurs fixes de localisation (vecteurs de dimension 306), avec sa file post-informatique stricte de contrôle anti-monopoles et contraintes budgétaires, prouve l'importance d'une algébrisation de proximité. Les batteries d'évaluation, cumulant tant la simulation synthétique de visites virtuelles (500 instances confirmant la convergence prédictive) et l'analyse fonctionnelle sémantique des requêtes RAG extrêmes bilingues ou de nature abstractive, légitiment techniquement la promesse originelle actée avec certitude sous des métriques de temps de réactions imperceptibles limitées à 100 ms.

Malgré tout le spectre abordé par ce programme académique ingénieux, ce sous-système connexe expose néanmoins par essence des latitudes et des limites technologiques concrètes. Le besoin inévitable de maintenir les calculs de matrice au sein de la mémoire vive pour l'ajustement temps réel des listes pose la problématique récurrente d'amplification d'infrastructure (\textit{scalability}) en cas de réplications à clusters multiples. Également, si les LLM locaux de modérés coefficients à l'image du qwen2.5:7b garantissent une autonomie d'inférence sécuritaire, l'intensité cognitive complexe d'un appel API pour orchestrer un graphe entier soudoie massivement parfois le répondant face aux API distantes tierces de classe 70b.

Les perspectives prospectives affluant à l'issue de cet acheminement invitent de prime évidence au tracé à grande échelle expérimentale vis-à-vis d'utilisateurs touristiques directs (méthode d'A/B Testing en masse). D'un point de vue strict du modèle d'exploration, ajuster organiquement la matrice adaptative de concept lors des passages saisonniers violents ou inclure d'autres systèmes de filtrage neuronal collaboratif d'appoint adosseraient inévitablement les inférences au sommet de la fiabilité prédictive, consacrant à terme pleinement Guidni non seulement comme un projet estudiantin visionnaire vis-à-vis de l'ingénierie tunisienne, mais comme un assistant de planification ubiquitaire irréfutable dans sa discipline.